\documentclass[sigplan,screen,nonacm]{acmart}
% sigplan = two-column SIGPLAN format
% nonacm  = removes ACM copyright block (for draft/preprint)

\usepackage{amsmath}
\usepackage{amssymb}
\usepackage{algorithm}
\usepackage{algorithmic}
\usepackage{booktabs}
\usepackage{listings}
\usepackage{xcolor}
\usepackage{url}

% Z80 assembly listing style
\lstdefinelanguage{z80asm}{
  morekeywords={LD,ADD,SUB,AND,OR,XOR,CP,NEG,INC,DEC,PUSH,POP,
    CALL,RET,JP,JR,DJNZ,EX,EXX,SBC,ADC,RET,HALT,NOP,
    RET NC,RET C,RET Z,RET NZ,
    DB,EQU,ALIGN},
  morecomment=[l]{;},
  sensitive=false,
}

\lstdefinelanguage{nanz}{
  morekeywords={fun,fn,if,else,while,for,in,return,let,var,global,
    import,struct,enum,true,false,void},
  morekeywords=[2]{u8,u16,i8,i16,bool,ptr},
  morecomment=[l]{//},
  morestring=[b]",
  sensitive=true,
}

\lstset{
  basicstyle=\ttfamily\footnotesize,
  columns=fullflexible,
  keepspaces=true,
  frame=single,
  framesep=2pt,
  xleftmargin=4pt,
  xrightmargin=4pt,
  aboveskip=6pt,
  belowskip=4pt,
  numbers=none,
  breaklines=true,
  captionpos=b,
}

% Compact itemize/enumerate
\usepackage{enumitem}
\setlist{nosep,leftmargin=*}

\begin{document}

\title{Per-Function Calling Convention Optimization\\for Irregular Register Architectures}

\author{Anonymous}
\affiliation{\institution{Anonymous}}
\email{anonymous@example.com}

\begin{abstract}
Compilers for register-constrained architectures with irregular register
files---where registers have distinct semantic roles (accumulator, counter,
pointer)---typically use a single fixed calling convention for all functions.
We formalize \emph{Per-Function Calling Convention Optimization} (PFCCO):
given a program's call graph and a cost model over register classes, choose
a calling convention for each internal function to minimize total transfer
cost.  We show that PFCCO is solvable optimally in $O(n)$ time for acyclic
call graphs with bounded parameter counts, by greedy dynamic programming in
reverse topological order.  Our implementation in the Nanz compiler
targeting the Zilog Z80 produces code that is 1.3--17$\times$ smaller and
up to 17$\times$ faster than SDCC~4.2.0 (which uses a state-of-the-art
empirically-optimized global convention), with emergent optimizations
including zero-byte function aliases.  We validate architecture-independence
by applying the same algorithm, unmodified, to the MOS~6502---an
architecture with even more extreme register irregularity (3~registers,
asymmetric capabilities)---achieving correct results across a five-path
cross-compilation oracle.
\end{abstract}

\keywords{calling conventions, register allocation, irregular architectures,
  Z80, 6502, compiler optimization, interprocedural optimization}

\maketitle

%% ============================================================
\section{Introduction}
\label{sec:intro}

Register allocation is among the most studied problems in compiler
construction.  For \emph{intraprocedural} allocation---assigning physical
registers to virtual registers within a single function---optimal algorithms
exist for structured programs on irregular architectures~\cite{krause2013}.
For the \emph{interprocedural} dimension---how values are passed between
functions---the standard approach is a fixed calling convention: every
function receives its first argument in the same register, returns its
result in the same register, and so on.

On architectures with large, regular register files (ARM, x86-64, RISC-V),
the calling convention is primarily about partitioning registers into
caller-saved and callee-saved sets.  The choice of \emph{which} register
holds a parameter matters little because all general-purpose registers have
identical capabilities.

On \emph{irregular} architectures---microcontrollers like the Zilog Z80,
Intel 8080, Padauk, or SM83---registers have semantic roles enforced by the
instruction set.  The \textbf{accumulator}~(A) is the only register that can
be an ALU source/destination.  The \textbf{counter}~(B) is the only register
usable with \texttt{DJNZ}.  The \textbf{pointer pair}~(HL) is the only pair
usable for memory addressing.  Choosing the wrong register for a parameter
forces the compiler to insert \emph{adapter moves}---register-to-register
transfers at call sites or inside the function body.  On a Z80 with only
7~main registers, each unnecessary \texttt{LD} costs 4~T-states and 1~byte.

Krause~\cite{krause2022} addressed this by empirically searching for the
best \emph{single} calling convention per architecture, testing thousands of
configurations against standard benchmarks.  This improved SDCC's Z80
backend significantly---enough to justify breaking ABI compatibility.  But a
global convention is inherently a compromise: the best choice for a leaf
function doing arithmetic differs from the best choice for a function that
forwards arguments to a callee.

We observe that internal functions need not share a convention at all.  Each
function can have its own \emph{contract}---a mapping from parameter/return
slots to register classes---chosen to minimize total cost across the
program.

\paragraph{Contributions.}
\begin{enumerate}
\item \textbf{Formalization} of Per-Function Calling Convention Optimization
  (PFCCO) as a cost-minimization problem over register class assignments on
  a call graph~(\S\ref{sec:formulation}).
\item \textbf{An $O(n)$ algorithm} for bounded parameter counts on acyclic
  call graphs, with a proof of optimality for
  DAGs~(\S\ref{sec:algorithm}).
\item \textbf{A production implementation} in the Nanz compiler targeting
  Z80, demonstrating 1.3--17$\times$ code size reduction over SDCC~4.2.0 on
  representative examples~(\S\ref{sec:eval}).
\item \textbf{Identification of emergent optimizations}---notably zero-byte
  function aliases (\texttt{EQU} collapse) that arise naturally from cost
  minimization~(\S\ref{sec:eval:equ}).
\item \textbf{Cross-architecture validation} on the MOS~6502, demonstrating
  that the same PFCCO algorithm (zero changes to optimizer code) produces
  correct, optimized calling conventions on a fundamentally different register
  file~(\S\ref{sec:eval:6502}).
\end{enumerate}

%% ============================================================
\section{Background}
\label{sec:background}

\subsection{The Z80 Register File}

The Zilog Z80 (1976) has 14 8-bit registers organized into groups with
distinct capabilities:

\begin{table}[h]
\centering\small
\caption{Z80 register roles and key instructions.}
\label{tab:z80regs}
\begin{tabular}{@{}lll@{}}
\toprule
Register & Role & Key instructions \\
\midrule
A        & Accumulator   & \texttt{ADD}, \texttt{SUB}, \texttt{CP}, \texttt{NEG} \\
B        & Counter       & \texttt{DJNZ} \\
C, D, E  & General       & \texttt{LD r,r'} only \\
HL       & Pointer pair  & \texttt{LD A,(HL)}, \texttt{INC HL} \\
DE       & Index pair    & \texttt{ADD HL,DE}, \texttt{EX DE,HL} \\
BC       & Counter pair  & \texttt{ADD HL,BC}, \texttt{PUSH BC} \\
\bottomrule
\end{tabular}
\end{table}

Additionally, a \emph{shadow register bank}
(\texttt{A',B',C',D',E',H',L'}) is accessible via \texttt{EXX}, which
atomically swaps all three main pairs with their shadows.  We discuss
shadow registers in \S\ref{sec:future}.

The register file is \emph{irregular}: no two registers are
interchangeable.  Moving a value from C to A costs 4~T-states
(\texttt{LD~A,C}); moving from HL to DE costs 4T (\texttt{EX~DE,HL}).
These asymmetries mean that the \emph{choice} of register for a parameter
directly affects instruction selection and code quality.

\subsection{The 6502 Register File}
\label{sec:6502regs}

The MOS~6502 (1975) has only 3~programmer-visible registers, making it the
most register-constrained popular architecture:

\begin{table}[h]
\centering\small
\caption{6502 register roles and key instructions.}
\label{tab:6502regs}
\begin{tabular}{@{}lll@{}}
\toprule
Register & Role & Key instructions \\
\midrule
A        & Accumulator   & \texttt{ADC}, \texttt{SBC}, \texttt{AND}, \texttt{ORA}, \texttt{EOR}, \texttt{CMP} \\
X        & Index          & \texttt{LDX}, \texttt{STX}, \texttt{(zp,X)} indirect \\
Y        & Index          & \texttt{LDY}, \texttt{STY}, \texttt{(zp),Y} indirect-indexed \\
\bottomrule
\end{tabular}
\end{table}

The 6502 is \emph{more} irregular than the Z80: A~is the \emph{sole} ALU
operand (all arithmetic and logical operations implicitly read and write~A),
X~and~Y are not interchangeable (X~supports \texttt{(zp,X)} pre-indexed
addressing while Y~supports \texttt{(zp),Y} post-indexed addressing), and
there are no 16-bit register pairs---16-bit values must reside in zero-page
memory.  Transferring between X~and~Y requires two instructions through~A
(\texttt{TXA;TAY} or \texttt{TYA;TAX}, 4~cycles).

This extreme irregularity makes PFCCO \emph{more} impactful: choosing the
wrong register for a parameter on the 6502 is catastrophic, since there is no
``just use another general register'' fallback.  We use the 6502 as a
cross-architecture validation target in \S\ref{sec:eval:6502}.

\subsection{Fixed vs.\ Per-Function Conventions}

A \emph{calling convention} specifies which registers hold parameters and
return values.  SDCC~4.2.0 uses Krause's~\cite{krause2022}
empirically-optimized convention:

\begin{table}[h]
\centering\small
\caption{SDCC 4.2.0 Z80 calling convention.}
\label{tab:sdcc-cc}
\begin{tabular}{@{}lccc@{}}
\toprule
Slot    & 8-bit & 16-bit & Return \\
\midrule
Param~1 & A     & HL     & A / DE \\
Param~2 & L     & DE     & --- \\
Param~3+ & \multicolumn{2}{c}{stack} & --- \\
\bottomrule
\end{tabular}
\end{table}

A \emph{per-function convention} allows each internal function to use a
different register assignment.  External functions (library, ROM, interrupt
handlers) retain a fixed convention.

\subsection{The Nanz Compiler}

Nanz is a systems language targeting Z80-class processors.  Its pipeline is:
%
\begin{center}
\small
Nanz $\to$ Parser $\to$ HIR $\to$ \fbox{PFCCO} $\to$ MIR2 $\to$ PBQP~RA $\to$
\begin{cases}
  \text{Z80 codegen} \\
  \text{6502 codegen} \\
  \text{C99 codegen} \\
  \text{QBE codegen}
\end{cases}
\end{center}
%
PFCCO runs after HIR-to-MIR2 lowering (which assigns Layer~1 defaults) and
before the PBQP register allocator.  The same MIR2 module---with PFCCO
contracts already assigned---is emitted by architecture-specific code
generators.  Only the cost table and codegen change; the optimizer is shared.

\subsection{Register Classes}
\label{sec:classes}

\begin{table}[h]
\centering\small
\caption{Register classes for Z80 parameters.}
\label{tab:classes}
\begin{tabular}{@{}lll@{}}
\toprule
Class & Registers & Type affinity \\
\midrule
\textsc{Acc}     & A       & u8 (ALU operand) \\
\textsc{Counter} & B       & u8 (loop count) \\
\textsc{General} & C,D,E   & u8 (general) \\
\textsc{Pointer} & HL      & u16, ptr \\
\textsc{Index}   & DE      & u16 \\
\textsc{Pair}    & BC      & u16 \\
\bottomrule
\end{tabular}
\end{table}

\begin{table}[h]
\centering\small
\caption{Register classes for 6502 parameters.}
\label{tab:6502classes}
\begin{tabular}{@{}lll@{}}
\toprule
Class & Location & Type affinity \\
\midrule
\textsc{Acc6502}  & A              & u8 (sole ALU register) \\
\textsc{IdxX}     & X              & u8 (pre-indexed addressing) \\
\textsc{IdxY}     & Y              & u8 (post-indexed addressing) \\
\textsc{ZP8}      & Zero-page byte & u8 (3-cycle load) \\
\textsc{ZP16}     & Zero-page pair & u16, ptr (only option) \\
\bottomrule
\end{tabular}
\end{table}

\noindent The 6502 has 4~plausible classes for u8 (vs.\ 3 on Z80) but only
1~class for u16 (\textsc{ZP16})---there are no register pairs.  The PFCCO
search space per function remains $O(1)$-bounded.

%% ============================================================
\section{Problem Formulation}
\label{sec:formulation}

\begin{definition}[Contract]
A \emph{contract} for function $f$ with $k$ parameter slots and $r$ return
slots is a vector $C(f) = (c_1, \ldots, c_k, c_{r_1}, \ldots, c_{r_r})$
where each $c_i$ is a register class compatible with the slot's type.
\end{definition}

\begin{definition}[Adapter cost]
The \emph{adapter cost} of contract $C(f)$ is:
\begin{equation}
\text{adapt}(f, C) = \sum_{i} \text{move}(c_i, \text{nat}(f,i))
  + \sum_{j} \text{move}(\text{nat}_r(f,j), c_{r_j}) \cdot |\text{ret}(f)|
\end{equation}
where $\text{nat}(f,i)$ is the register class naturally required by $f$'s
body for the $i$-th parameter.
\end{definition}

\begin{definition}[Transfer cost]
The \emph{transfer cost} at call edge $e = (\text{caller}, \text{callee})$
is:
\begin{equation}
\text{xfer}(e, C) = \sum_{k} \text{move}(\text{arg}_k, C(\text{callee})_k)
\end{equation}
\end{definition}

\begin{definition}[PFCCO]
Given call graph $G=(F,E)$, register file $R$, and cost function
$\text{move}$, find $C: F \to \text{Contracts}$ minimizing:
\begin{equation}
\text{cost}(C) = \sum_{f \in F} \text{adapt}(f, C(f))
  + \sum_{e \in E} \text{xfer}(e, C) \cdot w(e)
\label{eq:total-cost}
\end{equation}
subject to: no two parameter slots of the same function map to the same
uniquely-forced physical register.
\end{definition}

\paragraph{The \textsc{NaturalClass} function.}
The bridge between intraprocedural instruction selection and interprocedural
ABI optimization:

\begin{algorithmic}[1]
\REQUIRE function $f$, parameter index $i$
\FOR{each instruction using param $i$}
  \IF{ALU op (\texttt{ADD/SUB/AND/OR/CP})}
    \RETURN \textsc{Acc}
  \ELSIF{memory access (\texttt{Load/Store})}
    \RETURN \textsc{Pointer}
  \ELSIF{loop counter (\texttt{DJNZ})}
    \RETURN \textsc{Counter}
  \ELSIF{forwarded as arg $k$ to callee $g$}
    \RETURN $C(g)_k$ \COMMENT{transitivity}
  \ENDIF
\ENDFOR
\RETURN \textsc{General}
\end{algorithmic}

The transitivity case is the crucial insight: if a parameter is forwarded to
a callee, its natural class is whatever the callee expects.  This creates
chains of aligned ABIs through the call graph.

%% ============================================================
\section{Algorithm}
\label{sec:algorithm}

\subsection{Greedy DP on Reverse Topological Order}

\begin{algorithm}[t]
\caption{PFCCO-Optimize}
\label{alg:pfcco}
\begin{algorithmic}[1]
\REQUIRE Module $M$
\STATE $G \gets \text{BuildCallGraph}(M)$
\STATE $\text{order} \gets \text{TopoSort}(G, \text{leaves first})$
\STATE $\mathit{CS} \gets \emptyset$
\FOR{$f$ in order}
  \IF{$f$ is extern or empty}
    \STATE $\mathit{CS}[f] \gets \text{defaults}(f)$
    \STATE \textbf{continue}
  \ENDIF
  \STATE $\text{best} \gets \text{defaults}(f)$
  \STATE $\text{bestCost} \gets \text{Cost}(f, \text{best}, G, \mathit{CS})$
  \FOR{$c \in \text{CartesianProduct}(\text{plausible}(f))$}
    \IF{$\text{HasConflict}(c)$} \textbf{continue} \ENDIF
    \STATE $\text{cost} \gets \text{Cost}(f, c, G, \mathit{CS})$
    \IF{$\text{cost} < \text{bestCost}$}
      \STATE $\text{bestCost} \gets \text{cost}$; $\text{best} \gets c$
    \ENDIF
  \ENDFOR
  \STATE $\mathit{CS}[f] \gets \text{best}$
\ENDFOR
\RETURN $\mathit{CS}$
\end{algorithmic}
\end{algorithm}

\subsection{Three-Layer ABI System}

The algorithm operates within a three-layer system:
\begin{description}[style=unboxed]
\item[Layer~1: Positional defaults.]  A heuristic assigns classes by
  position and type (first u8 $\to$ \textsc{Acc}, first u16 $\to$
  \textsc{Pointer}).
\item[Layer~2: User annotations.]  Hard constraints
  (\texttt{@z80\_a}, \texttt{@z80\_hl}) respected by the optimizer.
\item[Layer~3: Contract optimizer (PFCCO).]  Algorithm~\ref{alg:pfcco},
  searching for better assignments over the call graph.
\end{description}

\subsection{Complexity}

For a function with $P$ parameter slots, the candidate count is at most
$3^P$ (three classes per u8 slot), reduced by conflict filtering.  Cost
evaluation is $O(|f| + |\text{callsites}|)$.

For bounded $P$ (typically $P \le 4$ on Z80):
\begin{itemize}
\item Per function: $O(3^4 \cdot |f|) = O(|f|)$
\item Total: $O(n)$ where $n$ = total program size
\end{itemize}

\subsection{Optimality for DAGs}

\begin{theorem}
For acyclic call graphs with bounded parameter counts,
Algorithm~\ref{alg:pfcco} produces an optimal solution in $O(n)$ time.
\end{theorem}

\begin{proof}
Topological order (leaves first) ensures every callee is decided before its
callers.  When processing $f$, we evaluate all candidate assignments.  The
cost function is exact: adapter costs within $f$ are computed directly, and
transfer costs use callees' final contracts.  Since $f$'s choice affects
only (a)~$f$'s internal adapter cost and (b)~transfer costs on edges from
$f$ to its callees---not costs within other functions---the locally optimal
choice is globally optimal.  This is a standard bottom-up DP argument: on a
DAG, if each subproblem is solved optimally given its children's solutions,
the global solution is optimal.
\end{proof}

\paragraph{Recursive functions.}
Cycle members are appended to the topological order with intra-cycle
callees' contracts unknown.  The algorithm uses default contracts for
undecided cycle members, yielding a heuristic.  SCC decomposition with
fixed-point iteration would restore optimality at cost
$O(\text{iters} \times |SCC|)$.

%% ============================================================
\section{Implementation}
\label{sec:impl}

PFCCO is implemented in 455~lines of Go in the Nanz compiler's \texttt{mir2}
package.

\begin{table}[h]
\centering\small
\caption{Implementation components.}
\label{tab:impl}
\begin{tabular}{@{}llr@{}}
\toprule
Component & Function & Lines \\
\midrule
\texttt{OptimizeContracts}  & Main loop             & 32 \\
\texttt{candidateChoices}   & Cartesian $+$ filter  & 47 \\
\texttt{inferNaturalClass}  & Body analysis          & 43 \\
\texttt{edgeCost}           & Transfer cost          & 43 \\
\texttt{BuildCallGraph}     & CG construction        & 38 \\
\texttt{topoSort}           & Kahn's algorithm       & 54 \\
\bottomrule
\end{tabular}
\end{table}

\paragraph{Cost table.}
Move costs between classes (Z80 T-states):

\begin{table}[h]
\centering\small
\caption{Inter-class move costs (T-states).}
\label{tab:costs}
\begin{tabular}{@{}lccccc@{}}
\toprule
From $\backslash$ To & Acc & Ctr & Gen & Ptr & Idx \\
\midrule
\textsc{Acc} & 0 & 4 & 4 & --- & --- \\
\textsc{Counter} & 4 & 0 & 4 & --- & --- \\
\textsc{General} & 4 & 4 & 0$^\dagger$ & --- & --- \\
\textsc{Pointer} & --- & --- & --- & 0 & 4$^\ddagger$ \\
\textsc{Index} & --- & --- & --- & 4$^\ddagger$ & 0 \\
\bottomrule
\multicolumn{6}{@{}l}{\footnotesize $^\dagger$0 if same register, 4 if
  different. $^\ddagger$\texttt{EX~DE,HL}.}
\end{tabular}
\end{table}

\paragraph{6502 cost table.}
The 6502 cost table encodes asymmetric transfer costs (cycles):

\begin{table}[h]
\centering\small
\caption{6502 inter-class move costs (cycles).}
\label{tab:6502costs}
\begin{tabular}{@{}lccccc@{}}
\toprule
From $\backslash$ To & Acc & IdxX & IdxY & ZP8 & ZP16 \\
\midrule
\textsc{Acc}   & 0 & 2 & 2 & 3 & 6 \\
\textsc{IdxX}  & 2 & 0 & 4$^*$ & 4 & --- \\
\textsc{IdxY}  & 2 & 4$^*$ & 0 & 4 & --- \\
\textsc{ZP8}   & 3 & 3 & 3 & 6$^\dagger$ & --- \\
\textsc{ZP16}  & --- & --- & --- & --- & 0 \\
\bottomrule
\multicolumn{6}{@{}l}{\footnotesize $^*$Requires routing through A: \texttt{TXA;TAY} (4~cycles).}\\
\multicolumn{6}{@{}l}{\footnotesize $^\dagger$\texttt{LDA~zp1; STA~zp2} (6~cycles, routes through A).}
\end{tabular}
\end{table}

\noindent The 6502 cost table has the same structure as the Z80 table but
encodes different physical constraints.  The PFCCO algorithm consumes this
table without modification---the same \texttt{OptimizeContracts} function
operates on both.

\paragraph{Pipeline position.}
PFCCO runs after HIR$\to$MIR2 lowering and before PBQP register allocation.
This ordering is essential: PFCCO sets interface classes; the allocator maps
virtuals to physicals within each function respecting the contract.

\subsection{Cross-Architecture Validation Oracle}
\label{sec:impl:oracle}

To verify that PFCCO produces correct calling conventions across
architectures, we employ a \emph{round-trip oracle}: each test function is
compiled through multiple independent paths and results are compared:

\begin{enumerate}
\item \textbf{MIR2~VM} --- interpret MIR2 directly (reference)
\item \textbf{mir2c} --- compile MIR2 to C99, compile with \texttt{cc},
  execute natively
\item \textbf{QBE} --- compile MIR2 to QBE~IL, assemble with \texttt{qbe},
  link with \texttt{cc}
\item \textbf{QBE round-trip} --- MIR2 $\to$ QBE~IL $\to$ \texttt{qbe2mir2}
  parse $\to$ mir2c $\to$ \texttt{cc}
\item \textbf{6502} --- compile MIR2 to 6502 assembly, assemble, execute in
  cycle-accurate emulator
\end{enumerate}

All five paths must agree on every test case.  This provides a strong
correctness guarantee: a PFCCO-assigned calling convention is validated not
just by one backend but by independent compilation pipelines that share
only the MIR2 IR and the PFCCO contracts.

%% ============================================================
\section{Evaluation}
\label{sec:eval}

We compare Nanz with PFCCO against SDCC~4.2.0, which uses
Krause's~\cite{krause2022} empirically-optimized global convention---the
state of the art for Z80 ABI selection.

\subsection{\texttt{abs\_diff}: Leaf Function}

\begin{lstlisting}[language=nanz,caption={\texttt{abs\_diff} in Nanz.}]
fun abs_diff(a: u8, b: u8) -> u8 {
    if a < b { return b - a }
    return a - b
}
\end{lstlisting}

\noindent Nanz output (PFCCO: \texttt{A=a, C=b}):
\begin{lstlisting}[language=z80asm]
abs_diff:
    SUB C       ; 4T
    RET NC      ; 11T/5T
    NEG         ; 8T
    RET         ; 10T
; 4 bytes, 27T worst
\end{lstlisting}

\noindent SDCC output (global: \texttt{A=a, L=b}): 10~bytes, 34T worst.
SDCC's use of L for the second argument forces a save-to-C and redundant
recomputation.  PFCCO chose \textsc{General}(C) because \texttt{SUB~C} is a
single-byte instruction that sets carry for the conditional return.

\subsection{GCD: Loop-Carried Register Persistence}

\begin{lstlisting}[language=nanz,caption={GCD.}]
fun gcd(a: u8, b: u8) -> u8 {
    while a != b {
        if a > b { a = a - b }
        else     { b = b - a }
    }
    return a
}
\end{lstlisting}

SDCC's global ABI places \texttt{a} in~C (after saving from~A), requiring
\texttt{LD~A,C} every loop iteration.  PFCCO keeps \texttt{a} in~A
throughout: 14~bytes vs.\ 19~bytes, zero reload overhead.

\subsection{\texttt{forEach}: Iterator Fusion + \textsc{ClassCounter}}

\begin{lstlisting}[language=nanz,caption={Lambda iterator.}]
fun sum_chain(buf: ^u8, n: u8) -> u8 {
    var s: u8 = 0
    buf.forEach(|x| { s = s + x }, n)
    return s
}
\end{lstlisting}

PFCCO assigned \texttt{n} to \textsc{Counter}(B), enabling \texttt{DJNZ}.
Combined with iterator fusion (sequential pointer traversal via
\texttt{INC~HL} instead of indexed access), the result is 12~bytes,
$\sim$38T/element vs.\ SDCC's 29~bytes, $\sim$88T/element
(\textbf{2.3$\times$} speedup).

SDCC puts the third parameter on the stack, requiring IX-relative addressing
at 19T per access in the loop.

\subsection{Multi-Return and EQU Collapse}
\label{sec:eval:equ}

\begin{lstlisting}[language=nanz,caption={Multi-return.}]
fun minmax(a: u16, b: u16) -> (u16, u16) {
    if a <= b { return (a, b) }
    return (b, a)
}
fun min_of(a: u16, b: u16) -> u16 {
    let (lo, _) = minmax(a, b)
    return lo
}
\end{lstlisting}

\begin{table}[h]
\centering\small
\caption{Multi-return: Nanz vs.\ SDCC.}
\label{tab:multiret}
\begin{tabular}{@{}lrrrr@{}}
\toprule
Function & \multicolumn{2}{c}{Nanz} & \multicolumn{2}{c}{SDCC} \\
         & Bytes & T & Bytes & T \\
\midrule
\texttt{swap}    & \textbf{2} & \textbf{14} & 26 & 236 \\
\texttt{min\_of} & \textbf{3}$^*$ & \textbf{10} & $\sim$40 & $\sim$200 \\
\bottomrule
\multicolumn{5}{@{}l}{\footnotesize $^*$0 bytes with \texttt{EQU} collapse.}
\end{tabular}
\end{table}

\paragraph{EQU collapse.}
Because PFCCO aligned \texttt{min\_of}'s contract with \texttt{minmax}'s
(both: HL, DE $\to$ HL), the tail call requires no setup.  The assembler
replaces \texttt{min\_of:~JP~minmax} with \texttt{min\_of~EQU~minmax}---a
\textbf{zero-byte function}.  This is emergent: it was not explicitly
programmed but falls out of cost minimization.

C cannot express multi-return, forcing SDCC to use pointer out-parameters
with stack frames.  \texttt{swap} is 13$\times$ smaller and 17$\times$
faster.

\subsection{Fibonacci: Recursive ABI Alignment}

Nanz returns u16 in HL (matching \texttt{ADD~HL,DE}'s expectation).
SDCC's global convention returns in~DE, requiring \texttt{EX~DE,HL} after
every recursive call (4T$\times$2 per recursion level) plus two stack saves
per call (22T).

\subsection{Summary}

\begin{table}[h]
\centering\small
\caption{Evaluation summary: Nanz (PFCCO) vs.\ SDCC 4.2.0.}
\label{tab:summary}
\begin{tabular}{@{}lrrrc@{}}
\toprule
Example & Nanz & SDCC & Size & T-state \\
        & bytes & bytes & ratio & ratio \\
\midrule
\texttt{abs\_diff} (u8)  & 4  & 10 & 2.5$\times$ & 1.3$\times$ \\
\texttt{abs\_diff} (u16) & 10 & 16 & 1.6$\times$ & 1.2$\times$ \\
\texttt{gcd}             & 14 & 19 & 1.4$\times$ & $\sim$1.5$\times$ \\
\texttt{swap}            & 2  & 26 & \textbf{13$\times$}  & \textbf{17$\times$} \\
\texttt{min\_of}         & 3  & $\sim$40 & \textbf{13$\times$+} & \textbf{20$\times$+} \\
\texttt{forEach}         & 12 & 29 & 2.4$\times$ & \textbf{2.3$\times$} \\
\texttt{fib}             & 26 & 28 & 1.1$\times$ & $\sim$1.2$\times$ \\
\bottomrule
\end{tabular}
\end{table}

Per-function ABI's advantage scales with \emph{structural complexity}: more
functions, deeper call graphs, multi-return.  For isolated leaf functions,
SDCC's global ABI is near-optimal.  The largest wins come from features C
cannot express (multi-return, lambda fusion) combined with per-function ABI
that eliminates adapter overhead.

\subsection{Cross-Architecture: 6502 Results}
\label{sec:eval:6502}

To validate that PFCCO generalizes beyond the Z80, we applied the same
algorithm---with zero changes to the optimizer---to the MOS~6502, substituting
only the cost table (Table~\ref{tab:6502costs}) and code generator.

\paragraph{Round-trip oracle results.}
We tested four representative functions through all five oracle paths.
Every path agreed on every test case (20/20):

\begin{table}[h]
\centering\small
\caption{6502 round-trip oracle: all 5 paths agree on all test cases.}
\label{tab:6502oracle}
\begin{tabular}{@{}llrl@{}}
\toprule
Function & Test case & Result & PFCCO assignment \\
\midrule
\texttt{abs\_diff(10,3)} & all 5 paths & 7  & \texttt{a$\to$A, b$\to$X} \\
\texttt{abs\_diff(3,10)} & all 5 paths & 7  & \\
\texttt{fib(8)}          & all 5 paths & 21 & \texttt{n$\to$A} \\
\texttt{fib(10)}         & all 5 paths & 55 & \\
\texttt{clamp(0,5,10)}   & all 5 paths & 5  & \texttt{x$\to$A, lo$\to$Y, hi$\to$X} \\
\texttt{clamp(15,5,10)}  & all 5 paths & 10 & \\
\texttt{max3(2,3,1)}     & all 5 paths & 3  & \texttt{a$\to$A, b$\to$X, c$\to$Y} \\
\texttt{max3(100,200,150)} & all 5 paths & 200 & \\
\bottomrule
\end{tabular}
\end{table}

\paragraph{PFCCO contract analysis.}
PFCCO's choices on the 6502 reveal the same optimization principles as on Z80
but adapted to a different register topology:

\begin{itemize}
\item \textbf{\texttt{abs\_diff}}: PFCCO placed \texttt{b} in~X, enabling
  \texttt{STX~tmp; SBC~tmp} (5~cycles).  On Z80, PFCCO placed \texttt{b} in~C,
  enabling \texttt{SUB~C} (4T, 1~byte).  Same principle---put the subtrahend in
  the cheapest ALU-accessible location---different physical realization.
\item \textbf{\texttt{clamp}}: PFCCO used all three registers (A, X, Y),
  avoiding any zero-page spill.  The cc65 compiler would route the third
  parameter through its software stack ($\sim$20~cycles overhead).
\item \textbf{\texttt{fib}}: Loop-carried variables placed in zero-page;
  PFCCO kept the loop counter path short by assigning \texttt{n} to~A.
\end{itemize}

\paragraph{Comparison with cc65.}
The cc65 compiler (the standard C compiler for 6502) uses a fixed ABI: first
parameter in~A, all subsequent parameters on a software stack.  Software stack
access costs $\sim$20~cycles per parameter vs.\ 2--3~cycles for register
passing.  PFCCO eliminates this overhead entirely for internal functions by
assigning all u8 parameters to registers or zero-page.

\paragraph{Code sharing.}
Table~\ref{tab:6502sharing} shows the code shared between Z80 and 6502
backends.  The PFCCO optimizer, PBQP allocator, and all MIR2 analyses are
100\% shared---only the cost table and code generator change.

\begin{table}[h]
\centering\small
\caption{Code sharing between Z80 and 6502 backends.}
\label{tab:6502sharing}
\begin{tabular}{@{}lrl@{}}
\toprule
Component & LOC & Shared? \\
\midrule
MIR2 IR + types         & $\sim$900  & 100\% \\
PFCCO optimizer         & $\sim$455  & 100\% \\
PBQP register allocator & $\sim$430  & 100\% \\
Optimization passes     & $\sim$1070 & 100\% \\
Liveness analysis       & $\sim$300  & 100\% \\
\midrule
Z80 cost table + codegen  & $\sim$2000 & Z80 only \\
6502 cost table + codegen & $\sim$930  & 6502 only \\
\bottomrule
\end{tabular}
\end{table}

\noindent The shared core ($\sim$3{,}155~LOC) is $3.4\times$ the size of the
6502-specific code ($\sim$930~LOC), demonstrating that PFCCO's
architecture-independence is not merely theoretical but realized in a
production implementation.

%% ============================================================
\section{Related Work}
\label{sec:related}

\paragraph{Intraprocedural RA.}
Krause~\cite{krause2013} presents optimal register allocation via tree
decomposition for structured programs on irregular architectures,
implemented in SDCC.  PFCCO is complementary: it sets the calling convention
that the allocator then respects.

\paragraph{Global ABI search.}
Krause~\cite{krause2022} empirically evaluates thousands of conventions per
architecture, selecting the best single convention.  PFCCO strictly
generalizes this: any global ABI is a special case where all functions
receive the same assignment.

\paragraph{Link-time per-function CC.}
De~Bus et~al.~\cite{debus2004} and Caldwell \&
Chiba~\cite{caldwell2017} optimize conventions at link time for ARM.  Both
target regular RISC files where the problem is primarily caller/callee-save
partitioning, not register \emph{class} assignment.  Neither provides formal
optimality results.

\paragraph{Interprocedural RA for RISC.}
Wall~\cite{wall1986} proposes link-time register allocation for MIPS.
LLVM~IPRA~\cite{llvm-ipra} propagates clobbered-register sets to reduce
caller-save overhead.  Both target large, regular files where register
choice is immaterial.

\paragraph{PBQP.}
Scholz \& Eckstein~\cite{scholz2002} formulate register allocation as PBQP.
PFCCO's cost function has natural PBQP structure.  For Z80's small candidate
spaces ($\le$729), exhaustive enumeration suffices.

\paragraph{MSVC /GL.}
Microsoft's whole-program optimization promotes stack parameters to
registers for internal functions~\cite{msvc-gl}.  No formal analysis is
published.

\begin{table}[h]
\centering\small
\caption{Gap analysis: prior work vs.\ PFCCO.}
\label{tab:gap}
\begin{tabular}{@{}lllll@{}}
\toprule
Work & Scope & Reg.\ file & Optimal? & Multi-arch? \\
\midrule
Krause '13      & Intra     & Irregular & Yes        & No$^\ddagger$ \\
Krause '22      & Global    & Irregular & Best-of-$N$ & No$^\ddagger$ \\
De Bus '04      & Per-fn    & Regular   & Heuristic  & ARM only \\
LLVM IPRA       & Inter     & Regular   & Exact$^\dagger$ & Yes \\
Wall '86        & Inter     & Regular   & Heuristic  & MIPS only \\
\textbf{PFCCO}  & \textbf{Per-fn} & \textbf{Irregular} & \textbf{Optimal} & \textbf{Yes (Z80+6502)} \\
\bottomrule
\multicolumn{5}{@{}l}{\footnotesize $^\dagger$For save/restore only, not class assignment.}\\
\multicolumn{5}{@{}l}{\footnotesize $^\ddagger$SDCC supports multiple backends but ABI is per-architecture, not parameterized.}
\end{tabular}
\end{table}

%% ============================================================
\section{Future Work}
\label{sec:future}

\paragraph{Shadow register bank.}
The Z80's shadow bank (\texttt{A',B',\ldots,L'}) provides 7~additional
locations via \texttt{EXX}.  Our implementation uses shadows for 32-bit
values (\textsc{DWord}: HL+H'L') but these assignments are hardcoded by
type, not optimized by PFCCO.  Extending the optimizer to shadow classes
requires modeling the EXX atomicity constraint: switching one pair forces all
three to switch simultaneously.

\paragraph{Recursive optimization.}
SCC decomposition with fixed-point iteration would restore optimality for
mutually-recursive clusters.

\paragraph{Profile-guided weighting.}
Dynamic profiling would weight hot edges, shifting ABI choices for critical
paths.

\paragraph{Additional architectures.}
The 6502 validation confirms PFCCO's portability.  Natural next targets
include the Intel~8080 (subset of Z80), SM83 (Game Boy), and Padauk (3-bit
stack, extreme constraints).  Each requires only a cost table and code
generator---the optimizer is unchanged.

\paragraph{Benchmark suite.}
Evaluation against Krause's~\cite{krause2022} benchmarks (Whetstone,
Dhrystone, Coremark) would enable direct comparison of per-function vs.\
global ABI optimization.

%% ============================================================
\section{Conclusion}
\label{sec:conclusion}

We have formalized Per-Function Calling Convention Optimization as a
cost-minimization problem on call graphs with register class assignments,
proved its optimal solvability for acyclic call graphs in linear time, and
demonstrated its effectiveness in a production compiler targeting both the
Zilog~Z80 and MOS~6502.

The key insight is that on irregular register files, the calling convention
is not merely an interface specification but an optimization variable with
direct impact on instruction selection.  By treating it as such, the
compiler discovers optimizations---including zero-byte function
aliases---that no fixed convention can achieve.

Cross-architecture validation on the 6502 demonstrates that PFCCO is not a
Z80-specific trick but a general technique for irregular register
architectures.  The same algorithm, consuming only a different cost table,
produces correct and optimized calling conventions on an architecture with
fundamentally different constraints (3~registers vs.\ 7, no register pairs,
asymmetric index registers).  A five-path round-trip oracle provides strong
correctness guarantees across independent compilation pipelines.

%% ============================================================
\bibliographystyle{ACM-Reference-Format}

\begin{thebibliography}{9}

\bibitem{krause2013}
Philipp~Klaus Krause.
\newblock Optimal Register Allocation in Polynomial Time.
\newblock In \emph{Compiler Construction (CC)}, 2013.

\bibitem{krause2022}
Philipp~Klaus Krause.
\newblock Efficient Calling Conventions for Irregular Architectures.
\newblock \emph{arXiv:2112.01397v2}, 2022.

\bibitem{debus2004}
Bruno De~Bus, Bjorn De~Sutter, Ludo Van~Put, Dominique Chanet, and
  Koen De~Bosschere.
\newblock Link-time optimization of ARM executables.
\newblock 2004.

\bibitem{caldwell2017}
Andrew Caldwell and Shigeru Chiba.
\newblock Reducing calling convention overhead in object-oriented programs.
\newblock 2017.

\bibitem{wall1986}
David~W. Wall.
\newblock Global Register Allocation at Link Time.
\newblock In \emph{SIGPLAN Notices}, 1986.

\bibitem{scholz2002}
Bernhard Scholz and Erik Eckstein.
\newblock Register Allocation for Irregular Architectures.
\newblock In \emph{LCPC/SCOPES}, 2002.

\bibitem{llvm-ipra}
Vivek~Pandya.
\newblock Interprocedural Register Allocation for LLVM.
\newblock GSoC 2016.

\bibitem{msvc-gl}
Microsoft.
\newblock \texttt{/GL} (Whole Program Optimization).
\newblock MSDN Documentation.

\end{thebibliography}

\end{document}
